% UTF8 表示使用通用国际字符作为内码，ctexart 表示这是一篇中文论文
\documentclass[UTF8]{ctexart} 
% 通过调用宏包 geometry 调整页面设置为 A4 纸，页边距 2.6 cm
\usepackage[a4paper,margin=2.6cm]{geometry}

\usepackage{amsmath, amssymb, amsthm} % 数学公式、符号、定理
\usepackage{graphicx}                 % 插图
\usepackage[colorlinks=true,linkcolor=blue]{hyperref} % 超链接与交叉引用

% 设置文章标题 (注：根据笔记内容，此为第九章支持向量机，保留了你的原格式)
\title{第八章讲义 (及第九章支持向量机)}
\author{}
\date{\today}

% 设置环境 theorem
\newtheorem{theorem}{定理}[section]
\newtheorem{lemma}[theorem]{引理}
\newtheorem{definition}[theorem]{定义}
\newtheorem{remark}{注记}[section]
\numberwithin{equation}{section}

% 自定义命令
\newcommand{\RR}{\mathbb{R}}
\newcommand{\CC}{\mathbb{C}}

\begin{document}

\maketitle

\section{Ch9 支持向量机 (Support Vector Machines)}

\subsection{9.1 最大间隔分类器 (Maximal Margin Classifier)}

\subsubsection{9.1.1 什么是超平面 (What Is a Hyperplane?)}

\begin{definition}
在 $p$ 维空间中，超平面（Hyperplane）是一个维度为 $p-1$ 的平坦仿射子空间 (flat affine subspace)。
\end{definition}

对于 $p$ 维特征空间，超平面的数学方程可以表示为：
\begin{equation}
\beta_0 + \beta_1 X_1 + \beta_2 X_2 + \dots + \beta_p X_p = 0
\end{equation}
为了书写简便，我们通常将其写成向量内积的形式。令特征向量 $X = (X_1, X_2, \dots, X_p)^T$，权重向量 $\beta = (\beta_1, \beta_2, \dots, \beta_p)^T$，则超平面方程等价于：
\begin{equation}
\beta^T X + \beta_0 = 0
\end{equation}

如果空间中的一个点 $X$ 满足 $\beta^T X + \beta_0 > 0$，则说明该点位于超平面的一侧；反之，若满足 $\beta^T X + \beta_0 < 0$，则该点位于超平面的另一侧。

\subsubsection{9.1.2 使用分离超平面进行分类 (Classification Using a Separating Hyperplane)}

假设我们有 $n$ 个训练观测样本：$x_1, \dots, x_n$，其中每个样本 $x_i \in \RR^p$。
对应有两类标签：$y_1, \dots, y_n \in \{-1, 1\}$。

如果存在一个超平面能够完美地将两类数据分离开来（Separating Hyperplane），那么它必须满足以下性质：
\begin{align}
\beta_0 + \beta_1 x_{i1} + \dots + \beta_p x_{ip} > 0 &\quad \text{如果 } y_i = 1 \\
\beta_0 + \beta_1 x_{i1} + \dots + \beta_p x_{ip} < 0 &\quad \text{如果 } y_i = -1
\end{align}
由于标签 $y_i$ 的符号与表达式的值的符号在分类正确时是一致的，上述两个条件可以被等价地统一为一个极其简洁的数学不等式：
\begin{equation}
y_i (\beta_0 + \beta_1 x_{i1} + \dots + \beta_p x_{ip}) > 0, \quad \forall i = 1, \dots, n
\end{equation}

\paragraph{测试集的分类规则：}
对于一个未知的测试观测点 $x^*$，我们只需将其坐标代入超平面函数 $f(x^*) = \beta_0 + \beta_1 x_1^* + \dots + \beta_p x_p^*$ 中，根据其符号（$\text{sgn}(f(x^*))$）来决定其所属的类别。

\begin{remark}[分类的置信度]
$|f(x^*)|$ 的大小（即量级）具有重要的几何意义。它表示点 $x^*$ 距离超平面的远近。$|f(x^*)|$ 的值越大，说明点离分类边界越远，我们对其分类结果的置信度（Confidence）就越高。
\end{remark}



\begin{proof}[证明：样本点到超平面的几何距离公式]


设超平面方程为 $H: \beta^T x + \beta_0 = 0$。对于空间中任意一点 $x_i$，设其在超平面 $H$ 上的正交投影点为 $x_i'$。

由于 $\beta$ 是超平面 $H$ 的法向量，因此投影向量 $\vec{x_i' x_i} = x_i - x_i'$ 必然与法向量 $\beta$ 平行。我们可以将其表示为：
\[ x_i - x_i' = d \cdot \frac{\beta}{\|\beta\|} \]
其中，常数 $d$ 的绝对值 $|d|$ 即为点 $x_i$ 到超平面的垂直几何距离。

等式两边同时左乘 $\beta^T$：
\[ \beta^T x_i - \beta^T x_i' = d \cdot \frac{\beta^T \beta}{\|\beta\|} \]
由于 $\beta^T \beta = \|\beta\|^2$，上式右侧可化简为 $d \|\beta\|$。

又因为投影点 $x_i'$ 位于超平面 $H$ 上，必须满足 $\beta^T x_i' + \beta_0 = 0$，即 $\beta^T x_i' = -\beta_0$。代入上式可得：
\[ \beta^T x_i + \beta_0 = d \|\beta\| \implies d = \frac{\beta^T x_i + \beta_0}{\|\beta\|} \]

取绝对值即得到点到超平面的几何距离：
\[ \text{Distance} = \frac{|\beta^T x_i + \beta_0|}{\|\beta\|} \]
考虑到对于完美分类的超平面，同号性质 $y_i(\beta^T x_i + \beta_0) > 0$ 恒成立，我们可以利用标签 $y_i \in \{-1, 1\}$ 将绝对值符号去掉，得到最终的几何距离公式：
\begin{equation}
 \text{Distance} = \frac{y_i(\beta^T x_i + \beta_0)}{\|\beta\|} \label{eq:distance_hyperplane_to_point}
\end{equation}
\end{proof}

\paragraph{支持向量 (Support Vectors)：}
在图 9.3 中，那些恰好落在间隔虚线边界上的点，被称为支持向量。
它们在几何上到最大间隔超平面的距离严格等于 $M$。它们之所以被称为“支持”向量，是因为它们“撑起”了整个间隔边界——\textbf{如果这些点发生移动，最大间隔超平面的位置也会随之改变}。相反，任何不在边界上的非支持向量，即使在合理范围内随意移动，也完全不会影响超平面的最终位置。这种只依赖少数关键点的性质，正是 SVM 的精妙之处。


\subsubsection{9.1.4 最大间隔分类器的构造 (Construction of the Maximal Margin Classifier)}

根据上述几何直观，寻找最大间隔超平面的过程，可以转化为一个带有约束的优化问题：

\begin{align}
\max_{\beta_0, \beta_1, \dots, \beta_p, M} \quad & M \\
\text{subject to } \quad & \sum_{j=1}^p \beta_j^2 = 1 \label{eq:norm_constraint} \\
& y_i (\beta_0 + \beta_1 x_{i1} + \dots + \beta_p x_{ip}) \ge M, \quad \forall i=1,\dots,n \label{eq:margin_constraint}
\end{align}

\begin{remark}[对约束条件的解析]
这里的约束条件 \eqref{eq:norm_constraint} 即 $\|\beta\| = 1$ 是极其关键的。
在解析几何中，点 $x_i$ 到超平面 $\beta^T x + \beta_0 = 0$ 的垂直几何距离公式为：
$$\text{Distance} = \frac{y_i (\beta^T x_i + \beta_0)}{\|\beta\|}$$
只有当我们强制限定法向量的长度 $\sum_{j=1}^p \beta_j^2 = 1$ 时，分母才变为 1。此时，不等式 \eqref{eq:margin_constraint} 左侧的代数表达式 $y_i f(x_i)$ 才严格等于真实的几何垂直距离。因此，整个优化问题的本质就是在所有能够正确分离数据的超平面中，最大化那个最小的几何距离 $M$。
\end{remark}


\paragraph{推导：从几何直观到凸二次规划的等价转化}

在结合几何直观时，我们得到的原始优化问题为：
\begin{align*}
\max_{\beta_0, \beta, M} \quad & M \\
\text{subject to } \quad & \|\beta\| = 1 \\
& y_i (\beta^T x_i + \beta_0) \ge M, \quad \forall i=1,\dots,n
\end{align*}
这个优化问题包含一个非凸的球面约束 $\|\beta\| = 1$（或者写成 $\sum \beta_j^2 = 1$），在数学优化中极其难以直接求解。为了将其转化为易于求解的凸优化问题，我们需要利用超平面的\textbf{尺度缩放不变性 (Scaling Invariance)}。

\begin{proof}[优化问题的等价转化推导]
\textbf{第一步：取消法向量的单位长度限制。}

回顾真正的几何距离公式为 $d = \frac{y_i (\beta^T x_i + \beta_0)}{\|\beta\|}$（公式\ref{eq:distance_hyperplane_to_point}）。如果不强制限制 $\|\beta\|=1$，我们的本质目标是最大化所有样本点中最小的几何距离：
\[ \max_{\beta, \beta_0} \left( \min_{i} \frac{y_i (\beta^T x_i + \beta_0)}{\|\beta\|} \right) \]

\textbf{第二步：利用尺度缩放不变性固定“函数间隔”。}

对于分离超平面方程 $\beta^T x + \beta_0 = 0$，如果我们将其参数 $\beta$ 和 $\beta_0$ 同时乘以任意一个正数常数 $k$，等式依然成立，超平面的几何位置完全不发生改变。
但是，点代入方程得到的值，即\textbf{函数间隔} $y_i (\beta^T x_i + \beta_0)$ 会相应地放大或缩小 $k$ 倍。

既然这种缩放不影响分类边界的几何本质，我们不妨主动选择一个特定的缩放比例，使得离超平面\textbf{最近}的样本点（即落在间隔边界上的支持向量），其函数间隔恰好等于 1：
\[ \min_{i} y_i (\beta^T x_i + \beta_0) = 1 \]
这就意味着，对于所有的训练样本，都必然满足：
\[ y_i (\beta^T x_i + \beta_0) \ge 1, \quad \forall i=1,\dots,n \]

\textbf{第三步：目标函数的等价转化。}

在上述特定的缩放标准下，代入原本的几何距离公式，最小的几何距离（即我们需要最大化的间隔 $M$）就简化成了：
\[ M = \frac{1}{\|\beta\|} \]

因此，“最大化间隔 $M$” 在数学上等价于“最大化 $\frac{1}{\|\beta\|}$”，进而等价于“最小化法向量的长度 $\|\beta\|$”。

为了在最优化过程中方便求导，并且符合现成的二次规划（Quadratic Programming, QP）求解器的标准形式，我们通常将其进一步等价替换为最小化 $\frac{1}{2}\|\beta\|^2$。

至此，原本带有非凸约束的复杂问题，被完美且等价地转化为了\textbf{标准的凸二次规划问题（支持向量机基本型）}：
\begin{align}
\min_{\beta_0, \beta} \quad & \frac{1}{2} \|\beta\|^2 \label{eq:svm_primal_obj} \\
\text{subject to } \quad & y_i (\beta^T x_i + \beta_0) \ge 1, \quad \forall i=1,\dots,n \label{eq:svm_primal_const}
\end{align}
\end{proof}


\begin{proof}[证明：最优超平面的拉格朗日对偶求解]
为了求解上述带有不等式约束的优化问题，我们为每个约束条件 \eqref{eq:svm_primal_const} 引入拉格朗日乘子 $\alpha_i \ge 0$，构造拉格朗日函数：
\[ L(\beta, \beta_0, \alpha) = \frac{1}{2} \beta^T \beta + \sum_{i=1}^n \alpha_i \left[ 1 - y_i(\beta^T x_i + \beta_0) \right] \]

根据 KKT 条件，最优解必然满足拉格朗日函数对原变量的偏导数为零：
\begin{enumerate}
    \item \textbf{对 $\beta$ 求偏导：}
    \[ \frac{\partial L}{\partial \beta} = \beta - \sum_{i=1}^n \alpha_i y_i x_i = 0 \implies \beta = \sum_{i=1}^n \alpha_i y_i x_i \]
    \textit{（注：这表明最优超平面的法向量仅仅是部分样本点的线性组合。）}
    
    \item \textbf{对 $\beta_0$ 求偏导：}
    \[ \frac{\partial L}{\partial \beta_0} = -\sum_{i=1}^n \alpha_i y_i = 0 \implies \sum_{i=1}^n \alpha_i y_i = 0 \]
\end{enumerate}

将上述两个结论代回拉格朗日函数中，展开并化简以消去原变量 $\beta$ 和 $\beta_0$：
\begin{align*}
L(\alpha) &= \frac{1}{2} \left( \sum_{i=1}^n \alpha_i y_i x_i \right)^T \left( \sum_{j=1}^n \alpha_j y_j x_j \right) + \sum_{i=1}^n \alpha_i - \sum_{i=1}^n \alpha_i y_i \left( \sum_{j=1}^n \alpha_j y_j x_j \right)^T x_i - \beta_0 \sum_{i=1}^n \alpha_i y_i \\
&= \frac{1}{2} \sum_{i=1}^n \sum_{j=1}^n \alpha_i \alpha_j y_i y_j x_i^T x_j + \sum_{i=1}^n \alpha_i - \sum_{i=1}^n \sum_{j=1}^n \alpha_i \alpha_j y_i y_j x_i^T x_j - 0 \\
&= \sum_{i=1}^n \alpha_i - \frac{1}{2} \sum_{i=1}^n \sum_{j=1}^n \alpha_i \alpha_j y_i y_j x_i^T x_j
\end{align*}

由此，我们将主问题转化为了仅依赖于拉格朗日乘子 $\alpha$ 的\textbf{对偶问题 (Dual Problem)}：
\begin{align*}
\max_{\alpha} \quad & \sum_{i=1}^n \alpha_i - \frac{1}{2} \sum_{i=1}^n \sum_{j=1}^n \alpha_i \alpha_j y_i y_j x_i^T x_j \\
\text{subject to } \quad & \alpha_i \ge 0, \quad i=1,\dots,n \\
& \sum_{i=1}^n \alpha_i y_i = 0
\end{align*}

通过 SMO 等算法解出极值 $\alpha^*$ 后，即可反求得法向量 $\beta = \sum_{i=1}^n \alpha_i y_i x_i$。
最后，利用 KKT 条件中的\textbf{互补松弛条件} $\alpha_i [1 - y_i(\beta^T x_i + \beta_0)] = 0$，任取一个 $\alpha_s > 0$ 的点（即支持向量），代入等式 $y_s(\beta^T x_s + \beta_0) = 1$ 即可求出截距 $\beta_0$：
\[ \beta_0 = y_s - \beta^T x_s \]
至此，最优分离超平面完全确定。
\end{proof}

\subsubsection{9.1.5 不可分的情况 (The Non-separable Case)}

如果两类数据点相互交叠，不存在任何一个可以将它们完美分离的超平面。在这种情况下，上述的优化问题将存在矛盾（无可行解），最大间隔分类器将会失效。

为了解决这一问题，我们需要对“完美分离”这一严苛条件进行放松，允许少数点跨越边界甚至被错误分类，这便引出了下一节的内容：\textbf{支持向量分类器 (Support Vector Classifier, 引入软间隔)}。

\end{document}